%% -*- coding:utf-8 -*-
\chapter{Satztypen und Expletiva}
\label{chap-expletives}

%\ili{Yiddish} besseres Beispiel nehmen, das wirklich V2 zeigt.

Kapitel~\ref{chap-verb-position}\is{Satztyp|(} hat die Analyse von Verberst- und Verbzweitsätzen behandelt. Dieses Kapitel befasst sich
mit eingebetteten Sätzen mit Komplementierern (\zb \emph{that} im \ili{Englisch}en) und eingebetteten Interrogativsätzen.

\section{Die Phänomene}

Dieser erste Abschnitt führt das Phänomen ein und besteht aus drei Teilen: Abschnitt~\ref{sec-complementizer} befasst sich mit eingebetteten
Sätzen, die durch einen Komplementierer eingeleitet werden, Abschnitt~\ref{sec-phen-interrogatives} beschreibt die
Struktur von Interrogativsätzen und Abschnitt~\ref{sec-phen-use-of-expletives} befasst sich mit Expletiva in V2- und Interrogativsätzen.

\subsection{Eingebettete Sätze, die durch einen Komplementierer eingeleitet werden}
\label{sec-complementizer}

Wie bereits erwähnt, sind \ili{Afrikaans}, \ili{Niederländisch} und \ili{Deutsch} SOV-Sprachen, und das zeigt sich auch in eingebetteten
Sätzen, die durch einen Komplementierer eingeleitet werden. (\mex{1}) ist ein Beispiel:

\ea Ich weiß, dass Aicke das Buch heute gelesen hat. \z

%% Stellung der anderen Konstituenten ist frei:
%% \eal
%% \ex Ich weiß, dass das Buch Max heute gelesen hat.
%% \ex Ich weiß, dass das Buch heute Max gelesen hat.
%% \zl

\noindent
Das \ili{Englisch}e, das eine SVO-Sprache ohne V2 ist, lässt nur SVO-Abfolge zu.
\ea[]{
I  know that Kim has read the book yesterday. \english
}
\z
Allerdings können Elemente in \emph{that}-Sätzen vorangestellt werden:
\eal
\ex[]{
I know that yesterday Peter came.
}
\ex[?]{
I know that bagels, he likes.
}
\zl

\noindent
Interessanterweise lässt das \ili{Dänisch}e, das ebenfalls eine SVO-Sprache ist, sowohl SVO-Abfolge (\mex{1}) als auch V2-Abfolge (\mex{2}) in Sätzen
zu, denen ein Komplementierer vorangeht:
\ea[]{
\label{ex-at-max-ikke-har-3}
\gll Jeg  ved, at   Gert ikke  har læst   bogen          {i dag}.\\
     ich weiß dass Gert nicht hat gelesen Buch.\textsc{def} heute\\\jambox{(\ili{Dänisch})}
\glt `Ich weiß, dass Gert das Buch heute nicht gelesen hat.'
}
\z
%(Negation hilft, Verbstellung zu bestimmen)

\eal
\ex[]{
\gll Jeg  ved, at   {i dag} har Gert ikke læst bogen.\\
     ich weiß dass heute    hat Gert nicht gelesen Buch.\textsc{def}\\\jambox{(\ili{Dänisch})}
}
\ex[]{
\gll Jeg  ved, at   bogen          har Gert ikke  læst {i dag}.\\
     ich weiß dass Buch.\textsc{def} hat Gert nicht gelesen heute\\
}
\zl
Das Beispiel in (\ref{ex-at-max-ikke-har-3}) enthält die Negation, um zu zeigen, dass wir es hier tatsächlich
mit der SVO-Abfolge zu tun haben. Ohne die Negation ist nicht klar, ob Nicht-V2-Sätze in
Sätzen zugelassen sind, die durch einen Komplementierer eingeleitet werden, da (\mex{1}a) das finite Verb in zweiter
Position hat. Wenn die Negation vorhanden ist, ist klar, dass wir einen V2-Satz vorliegen haben, wenn die Negation auf
das finite Verb folgt wie in (\mex{1}b), und dass wir keinen V2-Satz vorliegen haben, wenn das finite Verb auf die Negation folgt wie in
(\mex{-1}) und somit in dritter Position steht.
\eal
\settowidth\jamwidth{(V2 oder SVO)}
\ex
\gll at Gert har læst bogen\\
     dass Gert hat gelesen Buch.\textsc{def}\\\jambox{(V2 oder SVO)}
\glt `dass Gert das Buch gelesen hat'
\ex
\gll at   Gert har ikke læst bogen\\
     dass Gert hat nicht gelesen Buch.\textsc{def}\\\jambox{(V2)}
\zl
Bei komplementiererlosen Sätzen ist die V2-Abfolge die einzige mögliche Abfolge:
\eal
\settowidth\jamwidth{(V2 oder SVO)}
\ex[]{
\gll Gert har ikke læst bogen\\
     Gert hat nicht gelesen Buch.\textsc{def}\\\jambox{(V2)}
}
\ex[*]{
\gll Gert ikke har læst bogen\\
     Gert nicht hat gelesen Buch.\textsc{def}\\\jambox{(SVO)}
}
\zl

\noindent
Das \ili{Jiddisch}e und das \ili{Isländisch}e sind ebenfalls SVO-Sprachen. Die Sätze, die mit einem
Komplementierer kombiniert werden, sind V2:
\eal
\ex
\gll Ikh meyn  az   haynt hot Max geleyent dos bukh.\footnotemark\\
     ich denke dass heute hat Max gelesen das Buch\\\yiddish
\footnotetext{\citew[\page 58]{Diesing90a}}
\glt `Ich denke, dass Max das Buch heute gelesen hat.'

\ex% check!
\gll Ikh meyn  az   dos bukh hot Max geleyent.\\
     ich denke dass das Buch hat Max gelesen\\

\zl
\ea
\longexampleandlanguage{
\gll Engum         datt í hug,  að   vert  væri að reyna til     að kynnast honum.\footnotemark\\
     niemand.\DAT{} fiel zu Sinn dass wert wäre zu versuchen \PREP{} zu kennenlernen ihn\\}{Isländisch}
\footnotetext{\citew[\page 75]{Maling90a-u}}
\glt `Es kam niemandem in den Sinn, dass es sich lohnen würde zu versuchen, ihn kennenzulernen.'
\z

\subsection{Interrogativsätze}
\label{sec-phen-interrogatives}

Die OV-Sprachen bilden eingebettete Interrogativsätze, indem sie eine Phrase, die ein
Interrogativpronomen enthält\footnote{
Die meisten Interrogativpronomen beginnen im \ili{Deutsch}en mit \emph{w} und im \ili{Englisch}en mit \emph{wh}. Phrasen,
die ein Interrogativpronomen enthalten, werden \emph{w}"=Phrasen bzw. \emph{wh}"=Phrasen
genannt. Interrogativsätze werden manchmal als \emph{w}"=Sätze oder \emph{wh}"=Sätze bezeichnet.
}, einem
ansonsten SOV-Satz voranstellen. (\mex{1}) zeigt deutsche\il{Deutsch} Beispiele:
\eal
\ex Ich weiß, wer heute das Buch gelesen hat. \german
\ex Ich weiß, was Aicke heute gelesen hat.
\zl

\noindent
Da Sprachen wie das \ili{Deutsch}e Scrambling zulassen, könnten Sätze wie die in (\mex{0}) auch einfach auf
die Permutation der Argumente eines Kopfes zurückzuführen sein. Die Generalisierung über diese \emph{w}"=Sätze
ist jedoch, dass ein beliebiges \emph{w}"=Element vorangestellt werden kann. (\mex{1}) gibt ein \ili{deutsch}es Beispiel mit einer Fernabhängigkeit:
\ea
\label{ex-wissen-Vortrag-halen-nonlocal}
Ich weiß nicht, [über welches Thema]$_i$ sie versprochen hat,\\
{}[[einen Vortrag \_$_i$] zu halten].
\z


\noindent
Hier ist die Phrase \emph{über welches Thema} ein Argument von \emph{Vortrag},
das in die VP eingebettet ist, die \emph{zu halten} enthält, die wiederum unter
\emph{versprochen hat} eingebettet ist. Die Generalisierung über Interrogativsätze ist, dass ein
Interrogativsatz aus einer Interrogativphrase (\emph{über welches Thema}) und einem Satz besteht,
in dem diese Interrogativphrase irgendwo fehlt (\emph{sie versprochen
  hat, einen Vortrag zu halten}).

Im \ili{Deutsch}en ist die Abfolge der anderen Konstituenten frei wie in Aussagehauptsätzen und eingebetteten
Sätzen mit einem Komplementierer, die zuvor besprochen wurden.

\eal
\ex Ich weiß, was keiner diesem Eichhörnchen geben würde. \german
\ex Ich weiß, was diesem Eichhörnchen keiner geben würde.
\zl

\noindent
Das folgt aus dem bisher Gesagten, da Interrogativsätze einfach SOV-Sätze sind, aus denen eine
Konstituente extrahiert wurde. Die Möglichkeit, Konstituenten zu scramblen, wird durch das Extrahieren einer Phrase nicht beeinträchtigt.


Im \ili{Dänisch}en und im \ili{Englisch}en bestehen die Interrogativsätze ebenfalls aus einer Interrogativphrase und einem SVO-Satz,
in dem die Interrogativphrase fehlt:
\eal
\ex
\gll Gert har givet ham bogen.\\
     Gert hat gegeben ihm Buch.\textsc{def}\\\jambox{(\ili{Dänisch})}
\glt `Gert hat ihm das Buch gegeben.'
\ex
\gll Jeg ved, hvad$_i$ [Gert har givet ham \_$_i$].\\
     ich weiß was \spacebr{}Gert hat gegeben ihm\\
\glt `Ich weiß, was Gert ihm gegeben hat.'
\ex
\gll Jeg ved, hvem$_i$ [Gert har givet \_$_i$   bogen].\\
     ich weiß wem \spacebr{}Gert hat gegeben {} Buch.\textsc{def}\\
\glt `Ich weiß, wem Gert das Buch gegeben hat.'
\zl
\itdopt{Vielleicht hängt Extraktion ja doch mit Kasus zusammen und mit Passivierung.}
(\mex{0}a) zeigt den Satz mit SVO-Abfolge, (\mex{0}b) ist ein Beispiel mit dem sekundären Objekt als
Interrogativpronomen und (\mex{0}c) ist ein Beispiel mit dem primären Objekt als
Interrogativpronomen. Die Position, die die jeweiligen Objekte in nicht-interrogativen Sätzen wie (\mex{0}a)
haben, ist mit \_$_i$ markiert.

Das \ili{Jiddisch}e ist insofern besonders, als es auch in Interrogativsätzen V2-Abfolge hat \citep[Abschnitte~4.1, 4.2]{Diesing90a}: Interrogativsätze
bestehen aus einer Interrogativphrase, die aus einem V2-Satz extrahiert wird:

\ea
\gll Ir veyst efsher [avu            do    voynt Roznblat   der goldshmid]?\footnotemark\\
     Sie wissen vielleicht \spacebr{}wo da wohnt Roznblat der Goldschmied\\
\footnotetext{
\citew[\page 65]{Diesing90a}. Zitiert aus Olsvanger, \emph{Royte Pomerantsn}, 1949.
}
\glt `Wissen Sie vielleicht, wo Roznblat der Goldschmied wohnt?'
% The gollsing an translation differes in the original Rozenblatt
\z
%
%% \eal 
%% \ex
%% \label{vosmaks}
%% \gll Ikh veys nit   vos$_i$ [Max hot gegesn \_$_i$].\footnotemark\\
%%      I know not     what \hphantom{[}Max has eaten\\
%% \footnotetext{\citew[\page 68]{Diesing90a}.}
%% \glt `I do not know what Max has eaten.'

%% \ex%check
%% \gll Ikh veys nit   [vos                    hot Max gegesn].\footnotemark\\
%%      ich weiß nicht \hphantom{[}was heute hat Max gegessen\\
%% \footnotetext{\citew[68]{Diesing90a}.}
%% \glt `Ich weiß nicht, was Max heute gegessen hat.'
%\zl

\noindent
Die Variation ist also \emph{w}-Phrase + SOV, \emph{w}-Phrase + SVO und \emph{w}-Phrase + V2.\pagebreak

Zusätzlich zur Frage nach der Abfolge innerhalb des eingebetteten Satzes (SOV, SVO, V2) gibt es
Variation in der Finitheit des eingebetteten Satzes:
\eal
\label{ex-what-to-read}
\ex[]{
I wonder what to read.
}
\ex[]{
I wonder what I should read.
}
\ex[*]{
Ich frage mich, was zu lesen.
}
\ex[]{
Ich frage mich, was ich lesen soll / kann.
}
\zl
Das \ili{Englisch}e lässt Infinitive mit \emph{to} zu. Im Vergleich zu finiten Interrogativsätzen fügt die
infinitivische Form eine modale Bedeutung hinzu. Das \ili{Deutsch}e lässt keine nichtfiniten Interrogativsätze zu, wie
(\mex{0}c) illustriert.

\subsection{Die Verwendung von Expletiva zur Markierung des Satztyps}
\label{sec-phen-use-of-expletives}


Die germanischen\il{Germanisch} Sprachen verwenden die Konstituentenabfolge zur Kodierung des Satztyps: V2-Hauptsätze können
Aussagen oder Fragen sein, je nach dem Inhalt des präverbalen Materials und der
Intonation. Ebenso bestehen eingebettete Interrogativsätze aus einer \emph{w}"=Phrase und einem SVO-, SOV-
oder V2-Satz. Die Voranstellung einer Konstituente in einem V2-Satz geht mit bestimmten informationsstrukturellen
Effekten einher: Etwas ist das Topik oder der Fokus einer Äußerung. Für eingebettete Sätze ist es in
einigen Sprachen wichtig, dass die Struktur transparent ist, das heißt, dass wir \emph{w} + SVO- oder
\emph{w} + V2-Abfolge haben. Es gibt Situationen, in denen es unangebracht ist, ein Element voranzustellen, und in
solchen Situationen verwenden die germanischen\il{Germanisch} Sprachen zur Aufrechterhaltung einer
bestimmten Abfolge Expletiva, das heißt Pronomen, die keinen semantischen Beitrag leisten.

Das \ili{Deutsch}e verwendet das Expletivum \emph{es}, um die Position vor dem finiten Verb zu füllen, wenn keine andere
Konstituente vorangestellt werden soll.
\eal
\ex Drei Reiter ritten zum Tor hinaus.
\ex Es ritten drei Reiter zum Tor hinaus.
\zl

\noindent
Das \ili{Dänisch}e verwendet das Expletivum \emph{der}, um deutlich zu machen, dass eine Extraktion einer Konstituente stattgefunden hat
\citep[\page 169]{MOe2011a}:
\eal
\label{ex-danish-interrogative-expletive}
\ex[]{
\longexampleandlanguage{
\gll Politiet            ved   ikke,  hvem der           havde placeret bomben.\footnotemark\\
     Polizei.\textsc{def} weiß nicht wer \textsc{expl} hat platziert Bombe.\textsc{def}\\}{Dänisch}
\footnotetext{DK}%
\glt `Die Polizei weiß nicht, wer die Bombe platziert hat.'
}
\ex[*]{
\gll Politiet ved ikke,  hvem havde placeret bomben.\\
     Polizei.\textsc{def} weiß nicht wer hat platziert Bombe.\textsc{def}\\
}
\zl
Ohne das Expletivum hätte man ein Muster wie das in (\mex{0}b). In (\mex{0}b) haben wir die
normale SVO-Abfolge, und es ist für die Hörer*in nicht offensichtlich, dass das Muster aus einem extrahierten
Element (dem Subjekt) und einem SVO-Satz besteht, in dem es fehlt. Das ist transparenter, wenn ein
Expletivum in die Subjektposition eingefügt wird wie in (\mex{0}a). (\mex{1}) zeigt das anhand der
Analyse, die in Abschnitt~\ref{sec-analysis-expletives} vorgeschlagen wird: (\mex{1}a) zeigt die hypothetische
Struktur, die sich ergeben würde, wenn man annähme, dass das Subjekt \emph{hvem} `wer'
extrahiert wird. Es ergäbe sich eine sogenannte \emph{string vacuous movement}: Das Subjekt wird an eine Stelle
direkt neben seiner ursprünglichen Position bewegt. In (\mex{1}b) hingegen wird die Subjektposition vom Expletivum eingenommen, und
daher ist klar, dass der eingebettete Satz eine besondere Struktur hat. Das Expletivum \emph{der} ist ein overter Marker für
die Hörer*in oder Leser*in des Satzes, der ihn als eingebetteten Interrogativsatz markiert.
\eal
\ex[*]{
\gll [hvem$_i$     [\trace$_i$ havde placeret bomben]]\\
     \spacebr{}wer {} hat platziert Bombe.\textsc{def}\\\danish
}
\ex[]{
\gll [hvem$_i$     [der                    havde \trace$_i$ placeret bomben]]\\
     \spacebr{}wer \spacebr{}\textsc{expl} hat {} platziert Bombe.\textsc{def}\\
}
\zl

\noindent
Genauso verwendet das \ili{Jiddisch}e ein Expletivum in eingebetteten Interrogativsätzen (\emph{w} + V2), wenn
%the subject is extracted and
es kein anderes Element gibt, das informationsstrukturell für die präverbale Position geeignet ist.
%% \ea
%% \label{vosmaks}
%% \gll Ikh veys nit   [vos Max hot gegesn].\footnotemark\\
%%      ich weiß nicht \hphantom{[}was Max hat gegessen\\
%% \footnotetext{\citew[68]{Diesing90a}.}
%% \glt `Ich weiß nicht, was Max gegessen hat.'
%% \z
%% 
%% \item 
(\mex{1}) zeigt Beispiele aus \citet[\page 403--404]{Prince89a}:

\eal
\label{ex-Yiddish-interrogatives-expletive}
\ex[]{
\gll ikh hob  zi  gefregt ver es         iz beser  far ir.\\
     ich habe sie gefragt wer \textsc{expl} ist besser für sie\\\yiddish
\glt `Ich habe sie gefragt, wer besser für sie ist.'}
\ex[]{
\gll ikh hob  im  gefregt vemen es        kenen ale dayne khaverim.\\
     ich habe ihn gefragt wen \textsc{expl} kennen alle deine Freunde\\
\glt `Ich habe ihn gefragt, wen alle deine Freunde kennen.'}
\zl

(\mex{0}a) ist ein Beispiel mit einem Interrogativpronomen, das das Subjekt ist, und (\mex{0}b) ist ein
Beispiel, in dem die präverbale Position nicht von einem Argument von \emph{kenen} `kennen', sondern von
einem Expletivum gefüllt wird. Das Subjekt \emph{ale dayne khaverim} `alle deine Freunde' bleibt zurück und das Objekt
\emph{vemen} `wen' wird extrahiert, da es das Interrogativpronomen ist.

\if0
\subsection{Abhängige Sätze, die durch einen Komplementierer eingeleitet werden}

Die SOV-Sprachen (\ili{Afrikaans}, \ili{Niederländisch}, \ili{Deutsch}, \ldots) kombinieren einen Komplementierer mit einem Verbletztsatz.
\ea Ich weiß, dass Aicke das Buch heute gelesen hat. \z

Wie in V1- oder V2-Sätzen ist die Abfolge der Konstituenten im \mf in Sprachen wie dem \ili{Deutsch}en frei:
\eal
\ex Ich weiß, dass das Buch Aicke heute gelesen hat.
\ex Ich weiß, dass das Buch heute Aicke gelesen hat.
\zl

Im \ili{Englisch}en wird der Komplementierer mit einem Satz in SVO-Abfolge kombiniert.
\ea
I  know that Kim has read the book yesterday.
\z

%% 
%% Andere Stellungen sind nicht möglich:
%% \eal
%% \ex[*]{
%% I know that has Max read the book yesterday.
%% }
%% \ex[*]{
%% I know that yesterday Max has read the book.
%% }
%% \zl

Das \ili{Dänisch}e ist interessant, da es wie das \ili{Englisch}e SVO ist, aber zwei Möglichkeiten für Komplementierer"=
phrasen hat: \emph{at} `dass' kann mit einem Satz in SVO-Abfolge oder mit einem V2-Satz kombiniert werden:
\ea[]{
\label{ex-at-max-ikke-har-2}
\gll Jeg  ved, at   Gert ikke  har læst   bogen          {i dag}.\\
     ich weiß dass Gert nicht hat gelesen Buch.\textsc{def} heute\\\jambox{(\ili{Dänisch})}
\glt `Ich weiß, dass Gert das Buch heute nicht gelesen hat.'
}
\z

Die Negation \emph{ikke} kann verwendet werden, um die Verbstellung zu bestimmen. Da wir angenommen haben, dass adverbiales
Material an VPs adjungiert, wird angenommen, dass sich die Negation in (\mex{0}) mit \emph{har læst bogen i
  dag} `hat das Buch heute gelesen' kombiniert. Das finite Verb steht links von der Negation in (\mex{1}). Dies
passt gut zur Analyse der Verbstellung, die in Kapitel~\ref{chap-verb-position} eingeführt wurde: Das
finite Verb \emph{har} `hat' bildet eine Konstituente mit \emph{læst  bogen i dag}, an die sich die Negation anschließen kann, aber
aufgrund der Verbbewegung wird das Verb in die Initialposition gebracht und dann wird das Subjekt \emph{Gert} vorangestellt.
\ea[]{
\gll Jeg  ved, at   Gert har ikke  læst  bogen          {i dag}.\\
     ich  weiß dass Gert hat nicht gelesen Buch.\textsc{def} heute\\\jambox{(V2)}
}
\z

Wer nicht bereit ist zu glauben, dass Daten wie (\mex{0}) ein Beleg für eingebettete V2-Sätze sind, den werden die
folgenden Beispiele vielleicht noch mehr überzeugen:
\eal
\ex[]{
\gll Jeg  ved, at   {i dag} har Gert ikke læst bogen.\\
     ich  weiß dass heute   hat Gert nicht gelesen Buch.\textsc{def}\\\jambox{(\ili{Dänisch})}
}
\ex[]{
\gll Jeg  ved, at   bogen          har Gert ikke  læst {i dag}.\\
     ich weiß dass Buch.\textsc{def} hat Gert nicht gelesen heute\\
}
\zl
(\mex{0}a) zeigt einen Komplementierer + V2-Satz mit vorangestelltem \emph{i dag} `heute' und (\mex{0}b) zeigt
einen Komplementierer mit einem V2-Satz mit vorangestelltem \emph{bogen} `das Buch'. \emph{har} `hat' steht eindeutig
in zweiter Position und die erste Position ist von etwas besetzt, das nicht das Subjekt ist.

%% \ex[*]{
%% \gll Jeg  ved, at   {i dag} Gert har læst bogen.\\
%%      ich weiß dass heute   Gert hat gelesen Buch.{\sc def}\\
%% }
%%
%% \ex[*]{
%% \gll Jeg  ved, at   bogen          Gert har læst {i dag}.\\
%%      ich weiß dass Buch.{\sc def} Gert hat gelesen heute\\
%% }
%% \zl
%%

% Isländisch: Wikipedia (en)

\begin{itemize}
\item \ili{Jiddisch}: eingebettete Sätze sind V2 \citep[]{Diesing90a}:
\eal
\ex
\gll Ikh meyn  az   haynt hot Max geleyent dos bukh.\footnotemark\\
     ich   denke dass heute hat Max gelesen   das Buch\\
\footnotetext{\citew[\page 58]{Diesing90a}}
\glt `Ich denke, dass Max heute das Buch gelesen hat.'

\ex% check!
\gll Ikh meyn  az   dos bukh hot Max geleyent.\\
     ich denke dass das Buch hat Max gelesen\\

\zl

Isländisch:
\ea
\gll Engum         datt í hug,  að   vert  væri að reyna til     að kynnast honum.\footnotemark\\
     niemand.\DAT{} fiel zu Sinn dass wert wäre zu versuchen \PREP{} zu kennenlernen ihn\\%\icelandic
\footnotetext{\citew[\page 75]{Maling90a-u}.}
\glt `Es kam niemandem in den Sinn, dass es sich lohnen würde zu versuchen, ihn kennenzulernen.'
\z

\end{itemize}

\section{Interrogativsätze}

\begin{itemize}
\item Deutsch, Niederländisch, \ldots: \emph{w} + V-letzt:

\eal
\ex Ich weiß, wer heute das Buch gelesen hat.
\ex Ich weiß, was Aicke heute gelesen hat.
\zl

Interrogativnebensätze beginnen mit einer \emph{w}-Phrase.

\item Die \emph{w}-Phrase kann von weit her kommen:
\ea
Ich weiß nicht, [\gruen{über welches Thema}]$_i$ sie versprochen hat,\\
{}[[einen Vortrag \_$_i$] zu halten].
\z

\item Stellung der anderen Konstituenten ist frei:
\eal
\ex Ich weiß, was keiner diesem Eichhörnchen geben würde.
\ex Ich weiß, was diesem Eichhörnchen keiner geben würde.
\zl

\end{itemize}

\begin{itemize}
\item Dänisch: Interrogativnebensätze sind \emph{w} + SVO

\eal
\ex
\gll Jeg ved, hvad Gert har givet ham.\\
     ich weiß was Gert  hat gegeben ihm\\
\glt `Ich weiß, was Gert ihm gegeben hat.'
\ex
\gll Jeg ved, hvem Gert har givet   bogen.\\
     ich weiß wem  Gert hat gegeben Buch.{\sc def}\\
\glt `Ich weiß, wem Gert das Buch gegeben hat.'
\zl

\end{itemize}

\begin{itemize}
\item \ili{Jiddisch}: Interrogativnebensätze \emph{w} + V2 \citep[Abschnitte~4.1, 4.2]{Diesing90a}

%% \ea
%% %\ex
%% \label{vosmaks}
%% \gll Ikh veys nit   [vos Max hot gegesn].\footnotemark\\
%%      ich weiß nicht \hphantom{[}was Max hat gegessen\\
%% \footnotetext{\citew[68]{Diesing90a}.}
%% \glt `Ich weiß nicht, was Max gegessen hat.'

%% %% \ex%check
%% %% \gll Ikh veys nit   [vos              hot Max gegesn].\footnotemark\\
%% %%      ich weiß nicht \hphantom{[}was heute hat Max gegessen\\
%% %% \footnotetext{\citew[68]{Diesing90a}.}
%% %% \glt `Ich weiß nicht, was Max heute gegessen hat.'

\ea
\gll Ir veyst efsher [avu            do    voynt Roznblat   der goldshmid]?\footnotemark\\
     Sie wissen vielleicht  \spacebr{}wo da wohnt Roznblat der Goldschmied\\
\glt `Wissen Sie vielleicht, wo Roznblat der Goldschmied wohnt?' 
\footnotetext{
\citew[\page 65]{Diesing90a}. Quoted from Olsvanger, \emph{Royte Pomerantsn}, 1949.
}
\z
\end{itemize}

\fi

\section{Die Analyse}

Dieser Abschnitt befasst sich zunächst mit Sätzen, die durch einen Komplementierer eingeleitet werden
(Abschnitt~\ref{sec-embedded-cp}), und behandelt dann eingebettete Interrogativsätze in
Abschnitt~\ref{sec-interrogatives}. In Abschnitt~\ref{sec-analysis-expletives} wird eine Lexikonregel zur
Einführung von Expletiva vorgeschlagen, die in verschiedenen Sprachen eine Rolle bei der Markierung des Satztyps spielen.

\subsection{Eingebettete Sätze, die durch einen Komplementierer eingeleitet werden}
\label{sec-embedded-cp}

Die Analyse von Komplementiererphrasen im \ili{Afrikaans}, \ili{Niederländisch}en, \ili{Deutsch}en und \ili{Englisch}en ist unkompliziert: Der
Komplementierer wird mit einer nicht-invertierten verbalen Projektion kombiniert. Für die ersten drei Sprachen ist das ein
verb"=letzter Satz (SOV), und für das \ili{Englisch}e ist es ein SVO-Satz. Die jeweiligen Analysen sind in
Abbildung~\ref{fig-german-cp} und~\ref{fig-english-cp} angegeben.
\begin{figure}
\begin{forest}
sm edges
[CP
       [C [dass] ]
       [S
        [{NP[\type{nom}]} [niemand] ]
        [V$'$
          [{NP[\type{acc}]} [ihn] ]
          [V [kennt] ]
           ] ] ]
\end{forest}
\caption{Analyse der deutschen Komplementiererphrase als C + SOV}\label{fig-german-cp}
\end{figure}
\begin{figure}
\begin{forest}
sm edges
[CP
       [C [that] ]
       [S
        [{NP[\type{nom}]} [nobody] ]
        [VP
          [V  [knows] ]
          [{NP[\type{acc}]} [him] ] ] ] ]
\end{forest}
\caption{Analyse der englischen Komplementiererphrase als C + SVO}\label{fig-english-cp}
\end{figure}

\ili{Jiddisch}e Komplementierer selegieren einen V2-Satz. Die Analyse des Beispiels in (\mex{1}) ist in Abbildung~\ref{fig-yiddish-cp} dargestellt.
\ea
\gll Ikh meyn  az   haynt hot Max geleyent dos bukh.\footnotemark\\
     ich denke dass heute hat Max gelesen das Buch\\\yiddish
\footnotetext{\citew[\page 58]{Diesing90a}}
\glt `Ich denke, dass Max das Buch heute gelesen hat.'
\z
\begin{figure}
\begin{forest}
sm edges
[CP
        [C [az;dass] ]
        [S
          [{Adv$_i$} [haynt;heute] ]
          [{S/Adv}
            [{V \sliste{ S//V }}
              [V [hot$_k$;hat] ] ]
            [{S//\vSLASH Adv}
              [NP [Max;Max] ]
              [{VP//\vSLASH Adv}
                [{V//V}  [\_$_k$] ]
                [VP/Adv
                  [VP
                    [V [geleyent;gelesen] ]
                    [NP [dos bukh;das Buch,roof] ] ]
                  [Adv/Adv [\_$_i$] ] ] ] ] ] ] ]
%% \draw[connect] (Adv/Adv.north east)   [bend right] to (VP/Adv.south east);
%% \draw[connect] (VP/Adv.north east)    [bend right] to (V/\vSLASH Adv.south east);
%% \draw[connect] (V/\vSLASH Adv.north east)   [bend right] to (S//\vSLASH Adv.south east);
%% \draw[connect] (S//\vSLASH Adv.north east)  [bend right] to (S/Adv.east);
%% \draw[connect] (S/Adv.north east)     [bend right] to (Adv);
\end{forest}
\caption{Analyse der jiddischen Komplementiererphrase als C + V2}\label{fig-yiddish-cp}
\end{figure}
Die Analyse sieht kompliziert aus, aber es ist wirklich nur die Kombination eines Komplementierers mit einem V2-Satz.
Im V2-Satz ist \emph{haynt} `heute' extrahiert und das finite Hilfsverb \emph{hot} `hat'
ist in die V1-Position bewegt.

Die Unterschiede zwischen den Sprachen können dadurch erfasst werden, dass man den Komplementierer Sätze mit
unterschiedlichen Merkmal--Wert"=Kombinationen selegieren lässt. Während Komplementierer in SOV-Sprachen Verbletztprojektionen selegieren (\textsc{initial}$-$), selegiert das \ili{Jiddisch}e V2$+$-Sätze (das Ergebnis der Anwendung des
Füller-Kopf-Schemas) und das \ili{Dänisch}e spezifiziert keines dieser Merkmale am selegierten Satz. Da
im \ili{Dänisch}en nichts spezifiziert ist, kann der eingebettete Satz die Form haben, die der Rest der \ili{Dänisch}en
Grammatik zulässt: Er kann SVO oder V2 sein.

\subsection{Interrogativnebensätze}
\label{sec-interrogatives}

Wie die Datendiskussion gezeigt hat, wird die Phrase, die das Interrogativpronomen enthält, aus dem
übrigen Satz extrahiert. Die Voranstellung der \emph{w}"=Phrase ist wie die Voranstellung in V2-Sätzen. Im
topologischen Feldermodell wird die vorangestellte Phrase in deutschen\ili{Deutsch} Relativsätzen und Interrogativnebensätzen
dem \vf zugeordnet \citep[\page 48--49]{MuellerGT-Eng5}. Der Unterschied zwischen Interrogativnebensätzen und V2-Sätzen ist die Position des
Verbs: V2-Sätze haben das Verb in Initialposition, während es in
Interrogativ- und Relativsätzen in Endposition steht. Abbildung~\ref{fig-interrog-clause-simple} zeigt die Analyse des
Interrogativnebensatzes in (\mex{1}). Das Pronominaladverb \emph{worüber} wird aus
dem Rest des Satzes extrahiert.

\ea 
Ich weiß, [worüber]$_i$ [ \trace$_i$ sie spricht]. 
\z
\begin{figure}
\begin{forest}
sm edges
[S
  [PP [worüber]]
  [S/PP
    [PP/PP [\trace]]
    [V\rlap{$'$}
      [NP [sie]]
      [V [spricht]]]]]
\end{forest}
\caption{Analyse eines einfachen Interrogativnebensatzes}\label{fig-interrog-clause-simple}
\end{figure}

\noindent % since figure moves away
Abbildung~\ref{fig-interrog-clause-simple} verwendet die \slasch-Notation, die ich bisher verwendet habe. Um
komplexere \emph{w}"=Phrasen zu erfassen, werde ich denselben Trick wie für andere Fernabhängigkeiten
verwenden und Information über \emph{w}"=Pronomen an Mutterknoten weitergeben. Um das zu modellieren, werde ich
listenwertige Merkmale wie \spr und \comps verwenden. Abbildung~\ref{fig-interrog-clause-simple-slash} zeigt denselben Satz wie
Abbildung~\ref{fig-interrog-clause-simple}, aber mit zwei Merkmalen, die traditionell in der
Analyse von Fernabhängigkeiten verwendet werden \citep[Kapitel~4 und 5]{ps2}: \que und \slasch. Per Konvention werden die umrahmten Zahlen
vor XPs gesetzt, wenn die XP ein Argument ist, und sie folgen der XP, wenn die XP an einer
Fernabhängigkeit beteiligt ist. Der Grund dafür ist, dass unterschiedliche Teile der Information geteilt werden. Eine vollständige Erklärung
des Unterschieds erfordert ein tieferes Verständnis der Mechanismen und kann hier nicht gegeben werden. Die
interessierte Leser*in sei auf \citew[Kapitel~9]{MuellerLehrbuch4} oder auf \citew{Borsley:Crysman:2021a} verwiesen.
%has to wait until
%Chapter~\ref{chap-HPSG-light}.
PP\ibox{1}[\slasch \sliste{ \ibox{1} }] bedeutet, dass die relevante Information
über die PP in \slasch gestellt wird, das heißt in die Liste, die nach oben durchgereicht wird, bis ein passender
Füller gefunden wird (eine PP, deren relevante Eigenschaften mit dem Element in der
\slashl identifiziert werden können). Sobald ein Füller gefunden wurde, wird kein \slasch-Element mehr nach oben weitergegeben. Die \slashl des
obersten Knotens ist die leere Liste.

\begin{figure}
\begin{forest}
sm edges
[{S[\slasch \eliste]}
  [PP\ibox{1} [worüber]]
  [{S[\slasch \sliste{ \ibox{1} }]}
    [{PP\ibox{1}[\slasch \sliste{ \ibox{1} }]} [\trace]]
    [V\rlap{$'$}
      [NP [sie]]
      [V [spricht]]]]]
\end{forest}
\caption{Analyse eines einfachen Interrogativnebensatzes mithilfe des Merkmals \slasch}\label{fig-interrog-clause-simple-slash}
\end{figure}

Dieser Apparat kann auch für die Modellierung der Verhältnisse in der Interrogativphrase verwendet
werden. Abbildung~\ref{fig-interrog-clause-simple-que-slash} zeigt, wie die Information über das
Interrogativpronomen innerhalb der komplexen \emph{w}"=Phrase im Baum mithilfe des Merkmals \que nach oben weitergegeben wird.
\ea 
Ich weiß, [über welches Thema]$_i$ [ \trace$_i$ sie spricht]. 
\z
%
Abbildung~\ref{fig-interrog-clause-simple-que-slash} ist vollständig parallel zu
Abbildung~\ref{fig-interrog-clause-simple}, außer dass Information über das \emph{w}"=Wort
hinzugefügt ist. Der Inhalt von \que wird hier nicht angegeben, aber die \que-Liste von \emph{w}"=Wörtern enthält Information,
die für die Semantik benötigt wird: Das \emph{w}"=Wort zeigt an, wonach gefragt wird, und diese Information wird
bis zur Ebene des vollständigen Satzes nach oben weitergegeben (siehe \citealt{GSag2000a-u} zu Interrogativnebensätzen im Allgemeinen und insbesondere zu ihrer
Semantik).

\begin{figure}
\centerline{\begin{forest}
sm edges
[{S[\textsc{que} \eliste, \slasch \eliste]}
  [{PP\ibox{1}[\textsc{que} \sliste{ \ibox{2} }]} 
     [P [über]]
     [{NP[\textsc{que} \sliste{ \ibox{2} }]}
       [{Det[\textsc{que} \sliste{ \ibox{2} }]} [welches]]
       [N [Thema]]]]
  [{S[\slasch \sliste{ \ibox{1} }]}
    [{PP\ibox{1}[\slasch \sliste{ \ibox{1} }]} [\trace]]
    [V\rlap{$'$}
      [NP [sie]]
      [V [spricht]]]]]
\end{forest}}
\caption{Analyse eines einfachen Interrogativnebensatzes mithilfe von \slasch und \que}\label{fig-interrog-clause-simple-que-slash}
\end{figure}

Interrogativnebensätze werden durch eine spezielle Variante des Füller-Kopf-Schemas lizenziert, nämlich ein Schema,
das verlangt, dass die initiale Tochter (der Füller) etwas in ihrer \que-Liste hat (Abbildung~\ref{fig-interrogative-clause-schema}). Das hat zur Folge,
dass der Füller ein \emph{w}"=Wort enthalten muss.

\begin{figure}
\centerline{
\begin{forest}
[{H[\textsc{que} \eliste, \slasch \eliste]}
  [\ibox{1}{[\textsc{que} \sliste{ \etag{} }]}]
  [H{[\slasch \sliste{ \ibox{1} }]} ]]
\end{forest}
}
\caption{Beschränkungen des Interrogativnebensatz-Schemas, die für alle germanischen Sprachen gelten}\label{fig-interrogative-clause-schema}
\end{figure}

Zusätzlich zu den Spezifikationen in Abbildung~\ref{fig-interrogative-clause-schema} muss
festgehalten werden, ob das Verb initial steht oder nicht bzw.\ ob die rechte Tochter ein einfacher
SOV"=Satz wie im Deutschen\il{Deutsch} oder ein V2"=Satz wie im Jiddischen\il{Jiddisch} ist. Außerdem
müssen Beschränkungen hinsichtlich der Verbform spezifiziert werden. Das \ili{Deutsch}e lässt nur finite Verben zu, während
das \ili{Englisch}e finite Verben und Infinitive mit \emph{to} zulässt (siehe (\ref{ex-what-to-read}) auf S.\,\pageref{ex-what-to-read}).

Bisher haben wir deutsche\il{Deutsch} Beispiele mit vorangestellten Präpositionalobjekten
betrachtet. Abbildung~\ref{fig-wer-das-buch-liest} zeigt die Analyse von (\mex{1}) mit einem Subjekt als
\emph{w}"=Phrase:
\ea wer das Buch liest \z
\begin{figure}
\centerline{\begin{forest}
sm edges
[S
       [{NP[\snom]} [wer] ]
       [{S/NP[\snom]}
         [{NP[\snom]/NP[\snom]} [\trace] ]
         [V$'$
           [{NP[\sacc]}  [das Buch,roof] ]
           [V [liest] ] ] ] ]
\end{forest}}
\caption{Analyse eines deutschen Interrogativnebensatzes mit dem Subjekt als \emph{w}"=Wort}\label{fig-wer-das-buch-liest}
\end{figure}

Mit dem, was wir bisher haben, können wir also Interrogativnebensätze im \ili{Deutsch}en und in anderen SOV-Sprachen analysieren, aber
es gibt noch offene Fragen in Sprachen wie dem \ili{Dänisch}en, wo die Einfügung eines Expletivums in die Subjektposition
erforderlich ist, wenn nach einem Subjekt gefragt wird (siehe
(\ref{ex-danish-interrogative-expletive})). Auch das \ili{Jiddisch}e kann Expletivpronomen in die
präverbale Position in Interrogativnebensätzen einfügen (siehe
(\ref{ex-Yiddish-interrogatives-expletive}) auf S.\,\pageref{ex-Yiddish-interrogatives-expletive}). 
Expletivpronomina sind Gegenstand des nächsten Abschnitts.

\subsection{Eine Lexikonregel zur Einführung von Expletivpronomen}
\label{sec-analysis-expletives}

Die verschiedenen Typen von Expletivpronomen, die in Abschnitt~\ref{sec-phen-use-of-expletives} vorgestellt wurden, können -- vielleicht etwas überraschend -- durch
eine einfache Lexikonregel erfasst werden, die der \argstl eines Lexikonelements ein Expletivum hinzufügt
\citep[\page 180]{MOe2011a}:
\ea
\label{positional-expl-lr}
Lexikonregel zur Expletiveinführung\is{Lexikonregel!Expletiveinführung}:\\
\ms{
head & \ms[verb]{
       vform & fin \\
       }\\
arg-st & \ibox{1}\\
} $\mapsto$
\ms{
arg-st & \sliste{ NP[\type{lnom}]\upshape \sub{\type{expl}} } $\oplus$ \ibox{1}\\
}
\z
Die Anwendung der Lexikonregel ist auf finite Verben beschränkt, da positionelle Expletiva in
V2-Sätzen auftreten und diese immer finit sind. Expletiva in Interrogativnebensätzen werden verwendet, um
die Subjektposition zu füllen, um die Extraktion zu markieren, und das ist natürlich auch nur in finiten
Sätzen erforderlich.

Der Kasus des Expletivpronomens wird als lexikalischer Nominativ spezifiziert, was bedeutet, dass es
für die Kasuszuweisungsprinzipien unsichtbar ist. Der Nominativ wird der ersten NP mit strukturellem Kasus zugewiesen
(siehe S.\,\pageref{case-p}), und da das Expletivum lexikalischen Kasus hat, ändert sich an der
Kasuszuweisung an die anderen Elemente nichts. Dasselbe gilt
für strukturelle Akkusative: Die erste NP mit \emph{strukturellem} Kasus erhält Nominativ und alle anderen
Akkusativ. Das Expletivum beeinflusst die Kasuszuweisung nicht.

Ebenso ist die bisher vertretene Theorie der Kongruenz\is{Kongruenz|(} nicht betroffen: Das Verb kongruiert mit der
ersten NP mit strukturellem Kasus. Das macht die richtigen Vorhersagen für die Kongruenz im \ili{Isländisch}en, wo
das Verb mit Objekten im Nominativ kongruiert \citep[\page 460]{ZMT85a}.
\eal
\ex
\gll Hefur henni      alltaf þótt    Ólafur      leibinlegur?\footnotemark\\
     hat sie.\DAT{} immer gefunden Olaf.\NOM{} langweilig.\NOM{}\\\icelandic
\footnotetext{
\citew*[\page 451]{ZMT85a}
}
\glt `Hat sie Olaf immer langweilig gefunden?'
\ex
\label{ex-dat-subj-passive-ditransitive-icelandic-two}
\gll Konunginum voru gefnar ambáttir.\footnotemark\\
     der.König.\DAT{} wurden gegeben.\F.\PL{} Dienerinnen.\NOM.\F.\PL\\
\footnotetext{
\citew*[\page 460]{ZMT85a}
}
\glt `Dem König wurden Sklavinnen gegeben.'
\zl

\noindent
Und der Ansatz zur Kongruenz funktioniert auch für Fälle von Fernpassiv im \ili{Deutsch}en, wo das Subjekt nicht das erste Element
in einer \argstl ist. Eines der Beispiele in (\ref{bsp-auskosten-fernpassiv}) auf
S.\,\pageref{bsp-auskosten-fernpassiv} sei hier wiederholt:
\eal
\ex Sie erlauben uns nicht, den Erfolg auszukosten.
\ex\iw{auskosten} Der Erfolg wurde uns nicht auszukosten erlaubt.\footnote{
        \citew[\page 110]{Haider86c}%
}

\label{bsp-auskosten-fernpassiv-haider-zwei}
\zl
Siehe S.\,\pageref{ex-arg-st-fernpassiv-object-control} für die jeweiligen \argst-Listen.%
\is{Kongruenz|)}

Was die Position im Satz betrifft, so ist das Expletivum im \ili{Dänisch}en ein Subjekt. Das ist
genau das, was wir wollen, und was aus der allgemeinen Abbildung von \argst auf \spr und \comps in SVO-Sprachen
folgt. Die Analyse von (\mex{1}) ist in
Abbildung~\ref{fig-interrogative-Danish-subject-extraction} dargestellt.
\ea
\label{ex-hvem-der-læser-bogen}
\gll hvem der læser bogen\\
     wer \expl{} liest Buch.\textsc{def}\\\danish
\glt `wer das Buch liest'
\z
Der Lexikoneintrag für \emph{læser} `lesen' ist in (\mex{1}) angegeben:
\ea
Lexikoneintrag für \emph{læser} `lesen' mit Expletivsubjekt:
\ms{
spr    & \sliste{ NP[\type{lnom}]\upshape \sub{\type{expl}} }\smallskip\\
comps  & \sliste{ NP[\type{str}], NP[\type{str}] } \smallskip\\
arg-st & \sliste{ NP[\type{lnom}]\upshape \sub{\type{expl}}, NP[\type{str}], NP[\type{str}] } 
}
\z
Die Kasuszuweisungsprinzipien weisen der ersten NP mit strukturellem Kasus Nominativ und der zweiten
Akkusativ zu. Wie Abbildung~\ref{fig-interrogative-Danish-subject-extraction} zeigt, wird das Expletivsubjekt
als Spezifikator in der Subjektposition realisiert, und der Nominativ und der Akkusativ auf der
\compsl werden als Objekte realisiert. Das „Nominativobjekt“ wird extrahiert und als das
Interrogativpronomen realisiert.

\begin{figure}
\centerline{\begin{forest}
sm edges
[S
       [{NP[\snom]} [hvem;wer] ]
       [{S/NP[\snom]}
         [{NP[\lnom]} [der;\textsc{expl}] ]
         [{VP/NP[\snom]}
           [{\vbarSLASH NP[\snom]}
             [V [læser;liest] ]
             [{NP[\snom]/NP[\snom]} [\trace] ] ]
           [{NP[\sacc]} [bogen;Buch.\textsc{def} ] ] ] ] ]
\end{forest}}
\caption{Analyse von Interrogativnebensätzen im Dänischen mit Subjektextraktion}\label{fig-interrogative-Danish-subject-extraction}
\end{figure}

Während die Analyse für (\ref{ex-hvem-der-læser-bogen}) funktioniert, lässt sie auch Beispiele mit dem
Expletivsubjekt in eingebetteten nicht-interrogativen Sätzen zu, und diese sind ungrammatisch:
\eal
\ex[]{
\gll at   han ikke læser en bog\\
     dass er nicht liest ein Buch\\\danish
}
\ex[*]{
\label{ex-at-der-ikke-læser-han-en-bog}
\gll at   der     ikke læser han en bog\\
     dass \expl{} nicht liest er ein Buch\\
}
\zl
In der Formulierung der Grammatik muss gesagt werden, dass das Subjekt in den
Fällen extrahiert werden muss, in denen ein Expletivum als Subjekt verwendet wird. Das schließt (\mex{0}b) mit \emph{han}
`er' rechts vom finiten Verb in Objektposition aus, und es erzwingt die Extraktion des
Arguments mit strukturellem Nominativ wie in (\ref{ex-hvem-der-læser-bogen}). Die folgende
Beschränkung sorgt dafür:
\ea
Beschränkung über den Output der Lexikonregel zur Expletiveinführung\is{Lexikonregel!Expletiveinführung} für das \ili{Dänisch}e:\\
\ms{
arg-st & \sliste{ \_, NP\ibox{1}[\slasch \sliste{ \ibox{1} }] } $\oplus$ \etag\\
}
\z
Sie besagt, dass es ein erstes Element auf der \argstl geben muss (das Expletivum) und dass das darauf
folgende Element einen \slashw haben muss, der mit dem Element identisch ist. Die einzige Möglichkeit, diese
Anforderung zu erfüllen, besteht darin, das Verb mit einer Extraktionsspur zu kombinieren. Daher
ist (\ref{ex-at-der-ikke-læser-han-en-bog}) ausgeschlossen. Die leere Box in (\mex{0}) steht für einen beliebigen Wert. Sie kann
als leere Liste oder als Liste beliebiger anderer Länge instanziiert werden.

Wenden wir uns dem Jiddischen zu: Die Analyse von Interrogativnebensätzen im \ili{Jiddisch}en kann ein initiales Expletivum im V2-Satz
beinhalten, wenn die Sprecher*in das Subjekt oder ein anderes Element aus informationsstrukturellen Gründen für
diese Position als unangebracht empfindet. Abbildung~\ref{fig-Yiddish-interrogative-expletive} zeigt die Analyse von (\mex{1}):
\ea
\gll ver es leyent dos bukh\\
     wer \expl{} liest das Buch\\\yiddish
\glt `wer das Buch liest'
\z

\begin{figure}
\centerline{{\begin{forest}
sm edges
[S
       [{NP[\snom]} [ver$_i$;wer] ]
       [{S/NP[\snom]}
         [{NP[\lnom]} [es$_j$;\textsc{expl}] ]
         [{S/NP[\snom]/NP[\lnom]}
           [{V \sliste{S//V}}
             [V [leyent$_k$;liest] ] ]
           [{S//\vSLASH NP[\snom]/NP[\lnom]}
             [{NP[\lnom]/NP[\lnom]} [\trace$_j$] ] 
             [{VP//\vSLASH NP[\snom]}, s sep+=1em
               [\vbarDSLv 
                 [V//V [\trace$_k$] ]
                 [{NP[\snom]/NP[\snom]} [\trace$_i$] ] ]
               [{NP[\sacc]} [dos bukh;das Buch,roof] ] ] ] ] ] ]
\end{forest}}}
\caption{Analyse eines jiddischen Interrogativnebensatzes mit einem vorangestellten Expletivpronomen}\label{fig-Yiddish-interrogative-expletive}
\end{figure}
Die Analyse ist komplexer als die \ili{Dänisch}e, aber das liegt daran, dass das \ili{Jiddisch}e V2-Sätze
in Interrogativnebensätzen hat, die zusätzlich zur Extraktion eine Verbbewegung beinhalten. Interrogativnebensätze haben
zwei extrahierte Elemente: eines für V2, das Expletivum im Beispiel, und ein weiteres, das die
Interrogativphrase ist (\emph{ver} `wer' im Beispiel). Die Abbildung zeigt zwei Elemente nach einem /. In
der Notation mit dem Merkmal \slasch gäbe es eine Liste mit zwei Elementen.

Damit ist die Analyse finiter Interrogativnebensätze mit und ohne Expletivpronomen abgeschlossen,
aber zu Expletivpronomen im Allgemeinen ist noch mehr zu sagen. Ich habe oben eine Lexikonregel postuliert, die ein
Expletivelement hinzufügt, und das erfasst Expletivpronomen in \ili{Jiddisch}en Interrogativnebensätzen. Sie funktioniert nicht nur
für Interrogativnebensätze, sondern für V2 im Allgemeinen. \ili{Jiddisch}e deklarative V2-Sätze können ebenfalls initiale
Expletiva haben, ebenso wie deutsche\il{Deutsch} V2-Sätze. (\mex{1}a) ist ein deutsches\il{Deutsch} Beispiel. Wie (\mex{1}b) zeigt,
darf das Expletivum nicht im \mf erscheinen.

\eal
\ex[]{\label{ex-es-lachen-drei-Kinder}
Es lachen drei Kinder.
}
\ex[*]{
dass es drei Kinder lachen
}
\zl
In deutschen\il{Deutsch} deskriptiven Grammatiken wird dieses Expletivpronomen „Vorfeld-\emph{es}“ genannt, und es wird betont, dass es
nicht das Subjekt und kein Argument des Verbs ist \parencites[\page 129, 177, 371]{Eisenberg2004a}[§1263]{Duden2005-Authors}. Die Tatsache, dass das \emph{es} nicht im
\mf erscheinen kann, wird als Beleg für seine Nicht-Argumenteigenschaft gesehen. Wir haben jedoch gesehen, dass das Expletivum
im \ili{Dänisch}en in der Subjektposition realisiert wird, sodass die Idee, es über die germanischen\il{Germanisch} Sprachen hinweg
einheitlich als das initiale Element der \argstl zu behandeln, einen gewissen Reiz hat. Dennoch ist es
unbestreitbar, dass das \vf der einzige Platz ist, an dem dieses Expletivum im \ili{Deutsch}en und \ili{Jiddisch}en erscheinen kann.
%\itdopt{check Yiddish}
Das Problem kann gelöst werden, indem man die folgende Beschränkung zur Lexikonregel zur Expletiveinführung
in den Grammatiken des \ili{Deutsch}en und \ili{Jiddisch}en hinzufügt:\is{Voranstellung}
\ea
Beschränkung über den Output der Lexikonregel zur Expletiveinführung\is{Lexikonregel!Expletiveinführung} für das \ili{Deutsch}e und \ili{Jiddisch}e:\\
\ms{
arg-st & \sliste{ NP\ibox{1}[\slasch \sliste{ \ibox{1} }] } $\oplus$ \etag\\
}
\z
Diese Beschränkung besagt, dass das erste Element in der \argstl (das Expletivum) ein \slasch-Element
mit den relevanten Eigenschaften des Expletivums haben muss. Da das Expletivpronomen nichts
in \slasch hat, kann es nicht direkt mit den jeweiligen Lexikoneinträgen kombiniert werden. Die Spur hat
etwas in \slasch, das ist gerade ihr Wesen. In der Analyse von (\ref{ex-es-lachen-drei-Kinder}) kann
also eine Spur mit dem durch die Regel in (\ref{positional-expl-lr}) lizenzierten
Lexikonelement für \emph{lachen} kombiniert werden, und dann kann das Expletivum als Füller fungieren. Die Beschränkung in 
(\mex{0}) ist ähnlich wie die für das Dänische. Der Unterschied ist, dass das erste Element der \argstl
im Deutschen und Jiddischen extrahiert werden muss, während es im Dänischen das zweite Element ist.

Ein Problem bleibt: Die Extraktion von Expletiva muss satzbegrenzt sein:
\ea[*]{
Es$_i$ glaube ich, dass \_$_i$ drei Kinder lachen.
}
\z
Während die Extraktion prinzipiell Satzgrenzen überschreiten kann (siehe
(\ref{ex-wissen-Vortrag-halen-nonlocal}) auf S.\,\pageref{ex-wissen-Vortrag-halen-nonlocal}), ist das in
(\mex{0}) ausgeschlossen. Das ist jedoch kein spezielles Problem der vorliegenden Analyse des positionellen \emph{es},
sondern eine allgemeine Eigenschaft expletiver Elemente. (\mex{1}) zeigt ein Beispiel mit dem Wetter-\emph{es},
das eindeutig ein Argument des Verbs \emph{regnen} ist:
\ea[*]{
Es$_i$ glaube ich, dass \_$_i$ regnet.
}
\z
Was auch immer Beispiele wie (\mex{0}) ausschließt, erfasst also auch (\mex{-1}).

Schließlich bleibt ein Problem: Ich habe eine Lexikonregel angegeben, die Lexikonelemente für
Interrogativnebensätze mit einem Expletivum in der Subjektposition lizenziert, aber was noch fehlt,
ist eine Beschränkung, die Sätze ohne das Expletivum ausschließt. In der Grammatik gibt es bisher
nichts, was das tut. Es ist möglich, so etwas zu formulieren, aber die formalen Werkzeuge wurden in
diesem Buch nicht eingeführt. Die Leser*in sei für Details auf \citew[\page 185]{MOe2011a} verwiesen.

\section{Zusammenfassung}

Dieses Kapitel hat eine Analyse von abhängigen Sätzen, die durch einen Komplementierer eingeleitet werden, und von
Interrogativnebensätzen geliefert. Zusammen mit den V1- und V2-Sätzen, die in Kapitel~\ref{chap-verb-position} behandelt wurden, deckt dies
die wichtigsten Satztypen in den germanischen\il{Germanisch} Sprachen ab. Die Variation in diesen untergeordneten Sätzen
hängt mit dem zusammen, was wir zuvor gesehen haben: Die SOV-Sprachen haben SOV-Abfolge in eingebetteten Sätzen, und einige SVO-Sprachen
haben SVO-Abfolge, einige lassen sowohl SVO als auch V2 zu und einige lassen nur V2 zu. Interrogativnebensätze
enthalten einen Satz mit einer Lücke, und der Füller ist die Interrogativphrase, die ein
\emph{wh}-Wort im \ili{Englisch}en und ein entsprechendes Wort in den anderen germanischen\il{Germanisch} Sprachen enthält. Die \emph{wh}-Phrase
kann aus einem einzigen Interrogativpronomen bestehen oder intern komplex sein. Die Information
über das Interrogativpronomen muss aus semantischen und syntaktischen Gründen am obersten Knoten der Interrogativphrase
vorhanden sein. Der syntaktische Grund ist natürlich, dass man sicherstellen muss, dass
die vorangestellte Phrase überhaupt ein Interrogativpronomen enthält. Die Information wird vom
Interrogativpronomen durch denselben Mechanismus nach oben weitergegeben, der auch für die
Extraktion verwendet wird, nur wird für die Abhängigkeit in der Interrogativphrase das \que-Merkmal
verwendet. Das \ili{Dänisch}e und das \ili{Jiddisch}e nutzen Expletivpronomen für die
Sichtbarmachung der Struktur in Interrogativnebensätzen. Um
das zu erfassen, haben \citew{MOe2011a} eine Lexikonregel vorgeschlagen, die das Expletivum in die
\argstl einführt. Dieses Expletivum kann im \ili{Dänisch}en als Subjekt fungieren und im \ili{Deutsch}en und
\ili{Jiddisch}en als positionelles Expletivpronomen.%
\is{Satztyp|)}

\questions{
\begin{enumerate}
\item Wo werden Expletiva in den germanischen\il{Germanisch} Sprachen zur Markierung des Satztyps verwendet?
\item Ist der Satz, der auf den Komplementierer in (\mex{1}) folgt, ein SVO- oder ein V2-Satz?
\ea
\gll at   Gert har ikke læst bogen\\
     dass Gert hat nicht gelesen Buch.\textsc{def}\\\danish
\z
\end{enumerate}

}

\exercises{

\begin{enumerate}
\item Analysieren Sie die Interrogativnebensätze in (\mex{1}):
\eal
\ex Ich weiß, wen Kim kennt.
\ex
\gll Jeg ved, hvem der kende Kim.\\
     ich weiß wer \expl{} kennt Kim\\\danish
\glt `Ich weiß, wer Kim kennt.'

\zl

\item Analysieren Sie den Satz in (\mex{1}). Verwenden Sie Dreiecke für die NP und die PP.
\ea Es schwammen zwei Delphine neben dem Boot. \z
\end{enumerate}

}

%      <!-- Local IspellDict: de_DE -->
